\documentclass[12pt]{article}
\usepackage[utf8]{inputenc}
\usepackage[T1]{fontenc}
\usepackage{hyperref}
\title{Pa Language Resources for K-12 Curriculum Development}
\author{Compiled by Claude Sonnet 5.5 (claude-sonnet-5-5) for Andr\'as Kornai}
\date{}
\begin{document}
\maketitle
\section{Introduction}
Pa (ISO 639-3 ppt; also Pare, Akium-Pare) is a Papuan language spoken in Western (Fly) Province, Papua New Guinea, in the lowlands near the Elevala and Strickland rivers. It belongs to the Kamula--Elevala family, which some sources place within Trans-New Guinea. Speaker estimates in the sources found are about 5,900--6,000 (the task description says about 10,000). It should not be confused with the unrelated Pare language of Tanzania.
\section{Primary Educational Resources}
No primers, textbooks or school programmes for Pa were found in the search results.
\section{Linguistic and Reference Materials}
\begin{itemize}
\item Wikipedia, Pa language: \url{https://en.wikipedia.org/wiki/Pa_language}
\item New Guinea World (University of Zurich) entry for Pa, which reportedly lists sources including Voorhoeve's 1975 sketch and wordlists by Waldeck: \url{https://newguineaworld.linguistik.uzh.ch/families/trans-new-guinea/central-west-new-guinea/digul-river-ok/kamula-elevala-river/elevala-river/pa}
\item Wikipedia, Kamula--Elevala languages: \url{https://en.wikipedia.org/wiki/Kamula%E2%80%93Elevala_languages}
\end{itemize}
\section{Religious and Community Resources}
\begin{itemize}
\item Joshua Project people-group profile, Pa/Pare: \url{https://joshuaproject.net/people_groups/14243/PP}
\end{itemize}
\section{Digital and Audio Resources}
None found.
\section{Curriculum Development Context}
Documentation appears limited to a grammar sketch and wordlists. Any curriculum would first require orthography and materials development with the community; no community organisation was found in the searches.
\end{document}